\section{Related Work}
\label{S:related-works}

Our study connects the recognition of author-stated limitations and future work with rhetorical
sentence classification, LLM-supported scholarly-text processing, and analysis of reported research
directions. These areas share input documents and some methods, but their output claims differ.
We review that relationship before positioning the current extraction evaluation.

\subsection{Rhetorical annotation and scientific sentence classification}
\label{S:rw-background}

Argumentative zoning provides an early account of sentences according to their role in scientific
argument~\citep{teufel2002summarizing}. \citet{liakata2010corpora} describe AZ-II alongside Core
Scientific Concepts (CoreSC), illustrating two complementary views of scholarly content. Roles
such as future work, weaknesses in prior research, and failures of an author's approach can all
help locate research-relevant statements, but they do not constitute a single uniform author-gap
label. The attribution target and the distinction between a statement's rhetorical role and its
semantic content remain part of the annotation definition.

Full-text annotation is also an established setting. The Dr.~Inventor corpus contains 40 Computer
Graphics articles with multiple annotation layers, including argumentative roles and statements
about advantages or disadvantages~\citep{fisas2016multilayered}. The presence of limitation and
future-work information in separate layers is relevant to our task: it shows why a joint extraction
interface can draw on earlier annotation work without claiming to introduce the first full-text
resource containing both kinds of statement.

Abstract classification resources offer a different input condition. PubMed~200k~RCT
provides biomedical abstract roles, while CSAbstruct supports classification of computer-science
abstract sentences~\citep{dernoncourt2017pubmed,cohan2019pretrained}.
MAZEA and Emerald~110k broaden the available rhetorical annotation
schemes~\citep{dayrell2012rhetorical,stead2019emerald}. Such resources support document understanding,
but an abstract role such as conclusion or objective does not establish that a sentence expresses
an author-stated limitation. Treating every non-gap-specific role as a reliable negative would
also require checking whether gap statements can occur inside that role.

Context matters for classification because a short sentence may inherit its subject or purpose
from nearby text. \citet{jin2018hierarchical} use a hierarchical sequence model to incorporate
surrounding sentences in medical abstracts. \citet{brack2024sequential} examine cross-domain
multi-task learning for scientific sentence classification, and \citet{lan2024multilabel} study
multi-label classification with LLMs. These studies motivate explicit treatment of context, but
they do not establish that complete-sequence modeling is always appropriate for sparse gap labels
in full papers. \citet{zhang2024extracting} instead investigate context-enhanced problem/method
sentence extraction, an adjacent task in which local context helps interpret a target sentence.
Our verifier uses neighboring source sentences for attribution without replacing the task by
whole-document rhetorical classification.

\subsection{Author-stated limitation recognition}
\label{S:rw-authorstated}

\citet{kilicoglu2018automatic} recognize self-acknowledged limitations in clinical research using
rules and learned approaches. Their heading and lexical evidence provides a direct precedent for
our localization and cue stages. The important relationship is mechanistic: document structure
reduces the search space, while sentence-level evidence distinguishes limitation claims from
surrounding exposition. This is more specific than treating all methods that process scientific
text as equivalent extraction baselines.

\citet{lan2024automatic} extend clinical limitation work to 200 annotated RCT publications with
15 categories and 24 subcategories. Their reported binary detector $F_1$ of 0.821 differs from the
limitation-type classification result. Their at-scale application also characterizes limitation
reporting. The dataset is useful to us because it supplies author-gap labels and non-positive
sentences in an article-based setting. However, a reported candidate-pool recognition score cannot
be compared directly with a reconstructed full-stream score: section filtering, label prevalence,
training choices, and the recall denominator must first be aligned.

LimGen contributes author-written limitation sections from 4,068 ACL papers for suggestive
limitation generation~\citep{faizullah2024limgen}. These sections provide useful reference and
training material, but section membership is a form of supervision rather than proof that every
sentence is a semantic limitation. The distinction matters when a precision filter rejects a
method description, result statement, or checklist item occurring inside a Limitations section.
Nor is the original generation task an official binary-detection baseline simply because its
references can be transformed into sentence labels.

BAGELS combines source parsing and LLM processing with limitation generation across NLP, ML, and
biology~\citep{alazher2025bagels}. Its extraction evaluation reports different performance for
ACL and NeurIPS samples, which the authors discuss in relation to explicit versus scattered
limitation reporting. That comparison motivates evaluating structural accessibility, but does not
isolate section placement from every other corpus difference. Reviewer evidence is another
important distinction: author statements and reviewer-identified limitations are related sources
of critique, not the same attribution target. We preserve that distinction in comparisons with
our authors'-own extraction task.

\subsection{Future-work recognition and characterization}
\label{S:rw-future}

\citet{zhu2019recognizing} study future-research sentences in full JASIST articles using BERT-based
recognition and subsequent clustering. Their work links sentence extraction to thematic analysis:
retrieval locates source statements, whereas clustering characterizes what those statements concern.
The two stages therefore answer different questions, and a useful thematic grouping does not
independently verify that a stated research direction remains open.

ACL FWS-RC introduces an annotated future-work corpus and taxonomy~\citep{hao2020acl}.
The later recognition and classification study of \citet{zhang2023fws} compares machine-learning
methods and reports macro-$F_1=0.9073$ for recognition and weighted $F_1=0.7263$ for typing.
These numbers must remain attached to their own tasks and evaluation protocols. The resource is
particularly relevant to our future-work branch, while its curated sentence input makes it a
different test from parsing and scanning a complete article. We accordingly report the FWS
breadth experiment as sentence recognition rather than general full-PDF extraction.

\citet{suray2024future} analyze future-work statements in usable-security and privacy research and
investigate their relationship with citing publications. This extends the purpose of extraction
from locating sentences to examining how research agendas are articulated and subsequently used.
It also makes interpretation cautious: a future-work statement can be vague, and later research
can engage with a direction without repeating its wording. Such tracking would require additional
corpus and citation evidence beyond the stage-level recognition experiments reported here.

\subsection{LLM extraction, inferred gaps, and generated critique}
\label{S:rw-methods}

LLMs can serve several roles in scholarly-text analysis. An extractor selects source passages;
a verifier judges a candidate's fit to a definition; a generator proposes statements that need not
appear in the source. BAGELS and FutureGen include both extraction and generation workflows
\citep{alazher2025bagels,alazher2025futuregen}. FutureGen's retrieval-augmented generation of
future-work text is a downstream use of evidence, while our Stage~C returns only a decision about
an existing candidate. The shared use of an LLM does not make their outputs or evaluation measures
interchangeable.

GAPMAP distinguishes explicitly expressed knowledge gaps from gaps inferred through
context~\citep{salem2025gapmap}. Its evaluation includes multiple input units, from paragraphs
and sections to a small full-text pilot. This is relevant to the document-level interpretation
of gap claims, while also showing why dataset coverage must be separated from processing unit.
A collection of paragraph-level cases cannot by itself establish a full-article extractor's
localization performance, and inferred missing knowledge need not correspond to one quoted
sentence from the authors.

LimitGen is a separate benchmark from LimGen. It evaluates critical-limitation identification on
synthetic and reviewer-based AI research references~\citep{xu2025llms}. Tree-of-Concerns addresses
unstated limitations using hierarchical multi-agent critique~\citep{mishra2026tree}.
Both concern scientific weaknesses that a system may identify without the authors acknowledging
them. They therefore help delimit our objective and possible downstream extensions rather than
providing an automatically comparable extraction score. Retaining the original sentence can
support provenance, but it does not establish the substantive validity or novelty of a generated
criticism or research idea.

\subsection{Comparison dimensions and the research question addressed}
\label{S:rw-gap}

\Cref{T:related-work} compares the selected studies in chronological order by the
year of the cited version. It separates evaluated domain, target statement, origin of the reference,
method, dataset coverage, and processing unit. The gap-distribution column describes characterization
of extracted or annotated statements, not hypothesis tests of classifier performance. The final
column explains the connection of each study to our extraction design or research questions.
Dataset coverage denotes the evaluated resource rather than a claim about where the method can
be applied universally.

Several comparison decisions follow. Authors' own limitations and criticisms of prior work require
different attribution policies. Binary recognition and fine-grained typing require different
metrics and denominators. Full-article inputs and curated sentences impose different localization
burdens. Author-written, reviewer-written, and generated references have different evidential
origins. Finally, extraction effectiveness, corpus characterization, and research-idea usefulness
are separate outcomes. These distinctions motivate our dataset-specific tables and the separation
of input accounting from total system cost.

The present contribution concerns the allocation and evaluation of processing within an extractive
funnel. Structural filtering, lexical rules, learned classifiers, and LLM refinement all have prior
precedents. We examine their measured combination through a common full-article stage comparison,
auxiliary target-dependent recognition results, and input-volume analysis. The resulting comparison positions selective verification as an allocation strategy whose
effectiveness must be evaluated under an explicit target and input condition.

% Restores every study in the original supervisor-required comparison table.
\begin{landscape}
\begingroup
\setstretch{1.04}
\footnotesize
\setlength{\tabcolsep}{2.3pt}
\renewcommand{\arraystretch}{1.12}
\begin{longtable}{@{}>{\raggedright\arraybackslash}p{2.5cm}>{\raggedright\arraybackslash}p{2.3cm}>{\raggedright\arraybackslash}p{1.7cm}>{\raggedright\arraybackslash}p{2.2cm}>{\raggedright\arraybackslash}p{2.65cm}>{\raggedright\arraybackslash}p{2.65cm}>{\raggedright\arraybackslash}p{1.8cm}>{\raggedright\arraybackslash}p{2.6cm}>{\raggedright\arraybackslash}p{3.5cm}@{}}
\caption{Selected related work in ascending year order for the cited version. The relevance column explains each connection to the extraction task and research questions. Gap-distribution analysis excludes tests of classifier performance. The final row concerns the evidence reported here.}\label{T:related-work}\\
\toprule
\textbf{Study/year} & \textbf{Evaluated domain} & \textbf{Statement type} & \textbf{Statement origin} & \textbf{Approach} & \textbf{Dataset coverage} & \textbf{Processing unit} & \textbf{Gap-distribution analysis} & \textbf{Why related to our study}\\
\midrule
\endfirsthead
\multicolumn{9}{l}{\tablename\ \thetable\ (continued)}\\
\toprule
\textbf{Study/year} & \textbf{Evaluated domain} & \textbf{Statement type} & \textbf{Statement origin} & \textbf{Approach} & \textbf{Dataset coverage} & \textbf{Processing unit} & \textbf{Gap-distribution analysis} & \textbf{Why related to our study}\\
\midrule
\endhead
\midrule
\multicolumn{9}{r}{Continued on next page}\\
\endfoot
\bottomrule
\endlastfoot
2016: Dr. Inventor corpus \citep{fisas2016multilayered} & Computer Graphics & Limitations; future work & Author-stated; separate annotation layers & Human multi-layer annotation & 40 full articles & Sentences in full articles & Descriptive label distributions & Precedent for full-text annotations containing both statement types in distinct layers.\\
2018: Limitation recognition \citep{kilicoglu2018automatic} & Clinical research & Limitations & Author-stated & Rules; supervised and semi-supervised classification & 1{,}197 source articles & Sentences & Not investigated & Direct precedent for heading/keyword localization and learned author-limitation recognition (RQ1).\\
2019: Future-research recognition \citep{zhu2019recognizing} & Information science; JASIST & Future work & Author-stated & BERT classification; clustering & JASIST training and application collections & Sentences in full articles & Descriptive cluster analysis & Connects full-text future-work recognition to subsequent thematic analysis (RQ2).\\
2020: ACL FWS-RC \citep{hao2020acl} & NLP; ACL papers & Future work & Author-stated & Manual annotation; taxonomy & 4{,}024 papers; 3{,}067 future-work sentences & Sentences & Descriptive categories and yearly trends & Provides future-work labels/taxonomy; distinguishes recognition from characterization.\\
2022: FWS recognition/typing \citep{zhang2023fws} & NLP; academic articles & Future work & Author-stated & Lexical/neural recognition; type classification & Annotated FWS recognition/typing corpora & Sentences & Keyword/content and later-research similarity analysis & Direct recognition comparator; positive/macro/weighted F1 and sample protocols must align.\\
2024: LimGen \citep{faizullah2024limgen} & NLP; ACL Anthology & Limitations & Generated; author-stated references & Dataset; LLM generation & 4{,}068 papers & Paper & No corpus-level inferential analysis & Source of author-written limitation references; generation and section supervision differ from detection.\\
2024: Limitation categorization \citep{lan2024automatic} & Biomedical RCT publications & Limitations & Author-stated & PubMedBERT detection and type classification & 200 annotated articles; about 12k in application & Sentences & Corpus distributions; annotated-sample confidence intervals & Provides the principal author-limitation target; distinguishes detection, typing, and corpus analysis.\\
2024: LLM-SSC \citep{lan2024multilabel} & Biomedicine; computer science; physical/chemical sciences & Rhetorical roles & Author-written text; role labels & LLM prompting and fine-tuning & Multiple abstract datasets & Abstract sequence & Not investigated & Context-aware LLM sentence labeling is adjacent to candidate verification, with different labels.\\
2024: Problem/method extraction \citep{zhang2024extracting} & NLP & Problem; method & Author-written text; sentence labels & Context-enhanced transformer & 305 abstracts; 74 full papers & Target sentence with adjacent context & Not investigated & Motivates local context for sparse target sentences; problem/method labels are not gap labels.\\
2024: Cross-domain SSC \citep{brack2024sequential} & Multiple scientific domains & Rhetorical roles & Author text; rhetorical labels & Hierarchical cross-domain multi-task classification & Multiple abstract and full-paper datasets & Sentence sequence & Role and transfer analysis & Motivates detector/context alternatives and explicit transfer conditions (RQ2).\\
2024: Future work at SOUPS \citep{suray2024future} & Usable security/privacy & Future work & Author-stated & Source-statement and citing-publication analysis & SOUPS proceedings and citing publications & Statement; citing paper & Characterization and subsequent engagement & Motivates later corpus use; recognition alone does not establish a direction remains open.\\
2025: BAGELS \citep{alazher2025bagels} & NLP; ML; biology & Limitations & Author-stated; reviewer-identified; generated & Rules and parsing; LLM extraction and generation & ACL, NeurIPS, and PeerJ collections & Paper & Corpus-level analysis proposed as future work & Closest extraction/generation precedent; motivates source grounding and section-aware evaluation.\\
2025: FutureGen \citep{alazher2025futuregen} & NLP; ML & Future work & Author-stated; reviewer directions; generated & LLM extraction; RAG generation & Several thousand ACL/NeurIPS papers & Paper & No corpus-level inferential analysis & Future-work extraction feeds generation; motivates selective input while keeping the tasks distinct.\\
2025: GAPMAP \citep{salem2025gapmap} & Biomedicine; biodiversity; multi-domain pilot & Knowledge gaps & Explicit; context-inferred & Dataset; LLM gap detection and inference & About 1,500 mixed-unit inputs; small full-text pilot & Paragraph; section; paper & No corpus-level inferential analysis & Clarifies explicit versus inferred gaps and the importance of the input unit.\\
2025: LimitGen \citep{xu2025llms} & AI; NLP and ICLR papers & Limitations & Synthetic; reviewer-identified & Benchmark; LLM and agent evaluation & Synthetic and reviewer-based benchmark cases & Paper & Descriptive aspect distribution & Defines an adjacent critical-limitation task; reviewer/synthetic targets differ from quoted author claims.\\
2025: Sci-Sentence \citep{bolanos2025scisentence} & Literature reviews across disciplines & Gap; limitation; other roles & Human rhetorical labels; augmented training text & Seven-class models; LLM/encoder comparisons & 700 human-labelled sentences; augmented variant & Sentence & Rhetorical-class/error analysis & Auxiliary target includes prior-work criticism; requires label/prompt alignment (RQ2).\\
2026: Tree-of-Concerns \citep{mishra2026tree} & ML; NLP; computer vision & Limitations & Inferred; reviewer/citation references & Hierarchical multi-agent critique & ToC-Bench: 414 papers; 1,905 limitations & Paper & No corpus-level inferential analysis & Addresses unstated critique rather than author-stated extraction; bounds downstream claims.\\
\midrule
\textbf{This study} & Biomedical full articles; NLP future work; auxiliary literature gaps & Author limitations/future work; separate auxiliary gap task & Existing author text; dataset-specific labels & Structural localization; detectors; selective LLM verification & 40-article principal SAL stream; separate auxiliary subsets & Source sentence; article-level resampling & Not reported in the present extraction evaluation & Measures stage contribution (RQ1), target dependence (RQ2), and estimated input efficiency (RQ3).\\
\end{longtable}
\endgroup
\end{landscape}

